\section{Geometries guaranteeing local convergence}\label{art:positive}

The radial direction supplies the simplest comparison field. It is especially
natural for a homogeneous term: if $\gamma$ is finite, convex and positively
homogeneous of degree one, then
\[
 \gamma((1-\tau)r)=(1-\tau)\gamma(r),\qquad 0\le\tau<1.
\]
Its inward model is exact even when its angular derivatives jump. Adding a
smooth term leaves a quadratic inward modeling error. We make this observation
precise, prove its equivalence with radial secant control, and then use finite
polyhedral geometry to obtain the local theorem stated in the introduction.

\subsection{Inward models and radial secants}

Fix a minimizer $x_*$, write $F_*=F(x_*)$, and let
$K=\{x:F(x)\le b\}$ be compact, with $b>F_*$. For differentiable $x,z\in K$,
set
\[
 \begin{gathered}
 r_x=x-x_*,\qquad g=\nabla F(x),\qquad h=\nabla F(z),\\
 T=(h-g)^T(z-x),\qquad N=(h-g)^Tr_z.
 \end{gathered}
\]
The radial secant condition is the following requirement on positive bands:
\begin{equation}\label{art:radial-condition}
 \begin{gathered}
 \text{for every }\eta>0\text{ there is }C_\eta\ge0\text{ such that}\\
 N^2\le C_\eta T
 \quad\text{whenever }F_*+\eta\le F(z)\le F(x),\quad T>0.
 \end{gathered}
\end{equation}
Here all such differentiable endpoint pairs are included, whether or not they
are eligible for the fixed line search. The equivalent model condition is
\begin{equation}\label{art:inward-condition}
 \begin{gathered}
 \text{for every }\eta>0\text{ there are }M_\eta\ge0,
       \ \theta_\eta\in(0,1/2]\text{ such that}\\
 0\le F(x-\tau r_x)-F(x)+\tau\nabla F(x)^Tr_x
       \le\frac{M_\eta}{2}\tau^2,\quad 0\le\tau\le\theta_\eta,
 \end{gathered}
\end{equation}
at every differentiable $x\in K$ with $F(x)\ge F_*+\eta$.
The inward point need not be differentiable. It remains in $K$ by convexity.
The parameter $\tau$ is a fraction of the distance to $x_*$, rather than a
Euclidean step length.

\begin{theorem}[Radial equivalence]\label{art:radial-equivalence}
For finite convex $F$ and the compact sublevel $K$ above,
\eqref{art:radial-condition} and \eqref{art:inward-condition} are equivalent.
\end{theorem}
\begin{proof}
Choose $G,R>0$ bounding subgradient norms on $K$ and $\norm{x-x_*}$ there.
The subgradient bound was proved in Theorem~\ref{art:certificate}; supporting
inequalities then give $|F(x)-F(z)|\le G\norm{x-z}$ on $K$.
Use the supporting-plane error
\[
 D_F(z,x)=F(z)-F(x)-\nabla F(x)^T(z-x)\ge0.
\]
This is the usual Bregman expression \cite{BB}; only its supporting endpoint
needs to be differentiable. At differentiable endpoints,
$T=D_F(z,x)+D_F(x,z)$.

Suppose first that \eqref{art:inward-condition} holds, and fix a positive band,
with constants $M,\theta$. If $N\ge0$, compare the supporting plane at $x$
with the inward model at $z-\tau r_z$ to obtain
\[
 \tau N\le D_F(z,x)+\tfrac12M\tau^2.
\]
If $N<0$, reverse the endpoints and use the inward model at $x$:
\[
 \tau(T-N)\le D_F(x,z)+\tfrac12M\tau^2.
\]
Thus there are $A\in[|N|,2GR]$ and
$D\in\{D_F(z,x),D_F(x,z)\}$ such that
$\tau A\le D+M\tau^2/2$ for $0\le\tau\le\theta$.
For $M>0$ and $A\le M\theta$, choose $\tau=A/M$, giving $A^2\le2MD$.
If $A>M\theta$, take $\tau=\theta$ to obtain
$D\ge A\theta/2$ and hence $A^2\le(4GR/\theta)D$.
When $M=0$, the same endpoint choice gives an even stronger estimate.
Therefore
\begin{equation}\label{art:inward-to-radial}
 N^2\le \max\{2M,4GR/\theta\}\,T.
\end{equation}

Conversely, assume \eqref{art:radial-condition}, fix $\eta>0$, and put
\[
 \theta=\min\{1/2,\eta/(4GR)\}.
\]
For $x$ on the $\eta$-band and $z=x-\tau r_x$, $0<\tau\le\theta$,
convexity gives $F(z)\le F(x)$ and the Lipschitz estimate gives
$F(z)\ge F_*+3\eta/4$. Suppose initially that $z$ is differentiable and set
$a=(\nabla F(x)-\nabla F(z))^Tr_x\ge0$. Then
\[
 T=\tau a,\qquad N=-(1-\tau)a.
\]
If $a>0$, the secant bound on the $\eta/2$-band implies
$a\le C_{\eta/2}\tau/(1-\tau)^2$. If $a=0$, no division is needed.
In both cases convexity yields
\[
 0\le D_F(z,x)\le\tau a\le4C_{\eta/2}\tau^2.
\]

To remove differentiability at $z$, use the standard fact that a finite convex
function is differentiable almost everywhere \cite[Theorem~25.5]{Rockafellar70}.
For fixed $0<\tau<1$, the dilation $x\mapsto x-\tau(x-x_*)$ is invertible;
therefore points at which both the original and dilated points are
differentiable form a full-measure, hence dense, set. Approximate $x$ by such
points $x_j$ in the interior of $K$. This interior is nonempty since $b>F_*$,
and its closure is $K$: moving any point slightly toward $x_*$ gives value
strictly below $b$. The strict margin $3\eta/4>\eta/2$ ensures that the
approximating pairs lie on the required band for large $j$.

At a differentiability point $x$ of a finite convex function,
$\nabla F(x_j)\to\nabla F(x)$. Indeed, local subgradient bounds give compactness,
and every subsequential gradient limit satisfies the supporting inequality
at $x$, whose subgradient is unique. Passing to the limit in the preceding
function-value inequality now proves \eqref{art:inward-condition} with
$M_\eta=8C_{\eta/2}$. The case $\tau=0$ is immediate.
\end{proof}

\begin{corollary}[Radial convergence]\label{art:radial-convergence}
Under either equivalent condition, every admissible run starting in $K$
converges in value to $F_*$ and in distance to $\arg\min F$. With a unique
minimizer it converges to $x_*$.
\end{corollary}
\begin{proof}
Take $W_\eta(x)=r_xr_x^T$. Convexity gives
$g^Tr_x\ge F(x)-F_*\ge\eta$, proving visibility with $m_\eta=\eta^2$.
Equation~\eqref{art:radial-condition} gives secant control on eligible pairs.
Finally $r_z=r_x+s$ implies $\Pi_sr_z=\Pi_sr_x$, so transport holds with
$L_\eta=0$. Apply Theorem~\ref{art:certificate}.
\end{proof}

This condition controls only inward first-order modeling error. In a smooth
region its infinitesimal expression is $r_x^T\nabla^2F(x)r_x$, rather than the
largest Hessian eigenvalue. A bound almost everywhere on that expression alone
would miss a jump of the derivative along a ray; the function-value condition
\eqref{art:inward-condition} includes such jumps.

\subsection{Centered homogeneous terms and smooth perturbations}

\begin{theorem}[Centered smooth-plus-sublinear objectives]\label{art:centered}
Let $F:\R^n\to\R$ be finite and convex, and suppose $x_*$ minimizes $F$.
Let $U$ be an open convex set, let $\psi$ be convex and locally $C^{1,1}$ on
$U$, and let $\gamma:\R^n\to\R$ be finite, convex and positively homogeneous
of degree one. Suppose
\[
 F(x)=\psi(x)+\gamma(x-x_*)\qquad(x\in U).
\]
Every admissible run whose initial sublevel is compact and contained in $U$
satisfies $F(x_k)\to F(x_*)$. If $x_*$ is the unique minimizer, then
$x_k\to x_*$. Neither original summand is required to be minimized at $x_*$.
\end{theorem}
\begin{proof}
Let $K=\{F\le F(x_0)\}$. Since $g_0\ne0$, convexity implies
$F(x_0)>F(x_*)$. On compact convex $K\subset U$, local $C^{1,1}$ regularity
provides a uniform Lipschitz constant $L$ for $\nabla\psi$: take a finite
cover by neighborhoods carrying local constants and subdivide any segment in
$K$ into pieces in that cover. Summing the local bounds gives the maximum
of those constants along the segment. Integrating the gradient along an
inward segment consequently gives
\[
 \left|\psi(x-\tau r_x)-\psi(x)
             +\tau\nabla\psi(x)^Tr_x\right|
 \le\tfrac12L\tau^2\norm{r_x}^2.
\]
At a differentiable point of $F$ in $U$, the representation and smoothness of
$\psi$ imply differentiability of $\gamma$ at $r_x$. Differentiating its
homogeneity identity gives $\nabla\gamma(r_x)^Tr_x=\gamma(r_x)$.
Its contribution to the inward remainder is therefore exactly zero.
The upper bound in \eqref{art:inward-condition} follows with
$M=L\max_{x\in K}\norm{r_x}^2$; its lower bound follows from convexity of $F$.
Corollary~\ref{art:radial-convergence} proves the claim.
\end{proof}

For $F(u,v)=u^2+|v|$, take $\psi(u,v)=u^2$ and $\gamma(u,v)=|v|$.
Every sublevel is compact and the minimizer is unique, so every admissible
run converges to $(0,0)$. Guo--Lewis proves the weaker subsequential conclusion
for this example \cite[Corollary~4.4]{GL18}. Their Euclidean-norm convergence
theorem \cite[Corollary~4.2]{GL18} is an established special case also recovered
here. Lewis--Overton \cite[Section~5.1]{LO} analyzes the scalar absolute-value
model with its stated line-search conventions.
The proof above uses inward modeling directly. Convex-gradient cocoercivity
\cite{WW} supplies an alternative secant estimate, but no global
Lipschitz-gradient hypothesis on $F$ is required.

\subsection{Local finite partitions}

The distinction between a local objective hypothesis and a hypothesis on the
whole initial sublevel is handled by a simple convex trapping argument.

\begin{lemma}[A common local basin]\label{art:localization}
Let $F:\R^n\to\R$ be finite and convex with a unique minimizer $x_*$.
Every sublevel is compact. For each neighborhood $U$ of $x_*$, there is
$\varepsilon>0$ such that $\{F\le F(x_*)+\varepsilon\}\subset U$.
Thus any admissible run starting in that sublevel stays in $U$, independently
of its initial SPD matrix.
\end{lemma}
\begin{proof}
Choose $r>0$ with $\overline B_r(x_*)\subset U$.
Continuity, compactness of the sphere and uniqueness give
\[
 \omega=\min_{\norm{u}=r}\bigl(F(x_*+u)-F(x_*)\bigr)>0.
\]
For $\norm{x-x_*}\ge r$, convexity at the point at distance $r$ on the
segment from $x_*$ to $x$ gives
\[
 F(x)-F(x_*)\ge\frac{\norm{x-x_*}}r\omega.
\]
Consequently $F$ is coercive, hence has compact sublevels, and
$\varepsilon=\omega/2$ gives the asserted inclusion. Armijo preserves it.
\end{proof}

\begin{proof}[Proof of Theorem~\ref{art:polyhedral}]
First discard cells not containing $x_*$. Each such closed cell has an open
complement containing $x_*$; finiteness supplies one ball that avoids them all.
For every remaining cell, its extension domain $V_i$ contains $x_*$.
Shrink to a common ball $B_r(x_*)\subset U\cap\bigcap_{i:x_*\in P_i}V_i$,
and then shrink again so that the gradients of all remaining extensions have
one Lipschitz bound $L$ there. This uses only finitely many local extensions;
they need not have shared a domain before shrinking.

Let $x$ be a differentiability point of $F$ in this ball, and choose a cell
$P_i$ containing it. Convexity of $P_i$ places the whole inward segment
$x-\tau(x-x_*)$, $0\le\tau\le1$, in $P_i$. Both $F$ and $f_i$ agree
along that segment. Comparing their right derivatives at $\tau=0$ gives
\[
 \nabla F(x)^T(x-x_*)=\nabla f_i(x)^T(x-x_*).
\]
Only equality in the radial direction is needed, including at cell boundaries.
The smooth Taylor estimate along the same segment gives
\[
 0\le F(x-\tau r_x)-F(x)+\tau\nabla F(x)^Tr_x
       \le\tfrac12L\tau^2\norm{r_x}^2
       \le\tfrac12Lr^2\tau^2.
\]
The first inequality uses convexity of $F$, not convexity of the extension.
Lemma~\ref{art:localization} supplies a compact positive sublevel wholly
contained in this ball, chosen before the run or its matrix. The inward
condition holds on that sublevel, so Corollary~\ref{art:radial-convergence}
gives value and iterate convergence. Armijo gives the asserted containment.
\end{proof}

\begin{corollary}\label{art:plq}
A finite convex piecewise linear-quadratic objective with a unique minimizer
has the local convergence property of Theorem~\ref{art:polyhedral}.
The same holds for a convex locally $C^{1,1}$ loss plus a finite convex
piecewise linear-quadratic regularizer, provided the sum has a unique minimizer.
\end{corollary}
\begin{proof}
By piecewise linear-quadratic we mean a representation by quadratic polynomials
on a finite cover by closed full-dimensional polyhedral cells. Each polynomial is a smooth
extension of its cell formula; adding the smooth loss preserves the required
local $C^{1,1}$ extensions. Theorem~\ref{art:polyhedral} applies.
\end{proof}

This covers finite absolute-value and hinge penalties with locally $C^{1,1}$ convex losses,
under the stated uniqueness assumption. Different cells may have different
quadratic Hessians. A finite maximum of arbitrary quadratic formulas need not
admit such a polyhedral partition and is not automatically covered.
There is also no claim that iterates identify a cell or that rejected trials
remain in the local structural neighborhood.

For a further nonpolyhedral example, consider a locally $C^{1,1}$ convex loss plus finitely
many Euclidean residual norms $\norm{A_jx-b_j}$. Near a minimizer, a term with
nonzero residual is smooth, whereas a zero residual becomes
$\norm{A_j(x-x_*)}$, a centered sublinear term. Their sum has the local
representation in Theorem~\ref{art:centered}. Under uniqueness,
Lemma~\ref{art:localization} supplies the common low-starting basin.
The smoothness assertion here uses the Euclidean norm away from zero;
it does not apply to every spectral penalty at a nonzero rank-deficient point.
